\begin{center}
\section*{C}
\end{center}
\addcontentsline{toc}{section}{C}
\term{Candidatura}
Indica l'atto attraverso il quale uno studente o un gruppo di studenti mostra formalmente il proprio interesse alla partecipazione e allo sviluppo di un progetto proposto da un'azienda.
\term{Capitolato}
È un documento ufficiale che descrive in modo dettagliato le caratteristiche, le condizioni e i requisiti di un progetto fungendo da riferimento durante lo sviluppo del progetto.
\term{Caption}
La caption di una tabella o di una immagine è una breve descrizione che fornisce informazioni sul contenuto della tabello o della immagine stessa. Ha lo scopo di spiegare il contenuto visivo, mentre nel caso delle immagini può produrre una breve descrizione o un titolo per aiutare il lettore a comprendere il contesto della stessa.
\term{Casi d'uso}
Si tratta di una descrizione strutturata che indica come gli attori (utenti) interagiscano con il sistema, andando a rappresentare sequenze di azioni che mostrano come possa essere utilizzato il sistema.
\term{Consuntivo}
Il consuntivo è un rapporto che confronta le previsioni con i dati effettivi, al fine di valutare le variazioni e le differenze tra ciò che era previsto e ciò che è stato effettivamente raggiunto. Il consuntivo può essere utilizzato per analizzare il bilancio finanziario, il tempo impiegato, le risorse utilizzate e altre metriche di progetto. Fornisce una valutazione critica delle prestazioni e dell'efficacia del progetto, consentendo di prendere decisioni informate e di apportare eventuali correzioni necessarie per mantenere il progetto sul giusto percorso.
\term{Container}
I container sono una tecnologia che consente di raggruppare e isolare le applicazioni con l'intero ambiente di runtime. In questo modo è più semplice mantenere un comportamento e una funzionalità coerenti, spostando al contempo l'applicazione posta all'interno del container tra ambienti (sviluppo, test, produzione) e tra cloud pubblico, privato, ibrido e on premise. Essendo leggeri e portabili, i container offrono la possibilità di uno sviluppo più rapido per venire incontro alle esigenze aziendali quando emergono.
